\section{Torque as a coordinate: local charts and useful complements}
\label{sec:torque-coordinate}

Near a point at which \(\nabla_{\vect i}M\neq0\), neighbouring constant-torque
surfaces form a stack of nonintersecting sheets. Numbering those sheets by
their torque value is therefore a legitimate local coordinate choice. The
implicit-function and submersion theorems formalize this picture: two
additional independent functions \(\xi_2,\xi_3\) can be selected so that
\[
 \boldsymbol\xi=(M(\vect i),\xi_2(\vect i),\xi_3(\vect i))
\]
has a nonsingular Jacobian. In this chart, a curved surface \(M=M_0\) becomes
the coordinate plane \(\xi_1=M_0\). This statement is local; it neither proves
a globally single-valued inverse nor prescribes the best complements
\cite{agrachev2004}.

\begin{figure}[tb]
\centering
\begin{tikzpicture}[>=Latex,font=\small]
\begin{scope}[xshift=0cm]
 \draw[->] (-2,0)--(2.5,0) node[right]{\(i_d\)};
 \draw[->] (0,-1.6)--(0,2.1) node[above]{\(i_q\)};
 \foreach \k/\c in {-1.1/20,-.55/35,.55/55,1.1/70}
   \draw[Purple!\c,very thick,domain=.35:2.15,samples=50]
     plot ({\x},{\k/\x});
 \node[align=center] at (0,-2.05){curved constant-\(M\) sheets\\in a current slice};
\end{scope}
\draw[->,very thick] (3.0,.2)--(5.0,.2) node[midway,above]{\(\boldsymbol\Phi\)};
\begin{scope}[xshift=7.2cm]
 \draw[->] (-2,0)--(2.4,0) node[right]{\(\xi_2\)};
 \draw[->] (0,-1.6)--(0,2.1) node[above]{\(M\)};
 \foreach \y/\c in {-1.0/20,-.5/35,.5/55,1.0/70}
   \draw[Purple!\c,very thick] (-1.9,\y)--(1.9,\y);
 \node[align=center] at (0,-2.05){parallel \(M=M_0\) planes\\in a task slice};
\end{scope}
\end{tikzpicture}
\caption{A regular task chart straightens constant-torque surfaces. It does
not guarantee a unique global inverse.}
\label{fig:torque-surface-task-plane}
\end{figure}


Useful complements answer different questions:
\begin{center}
\begin{tabular}{lll}
\toprule
Chart & Meaning of complements & Principal risk\\
\midrule
\((M,\alpha,z_0)\) & factor sharing; loss-null allocation & \(M=0\) singularity\\
\((M,\psi_a,\sigma)\) & active flux; excitation sharing & branch in \(i_q=M/(k_T\psi_a)\)\\
\((M,\psi_s,i_{\mathrm{exc}})\) & voltage-relevant flux; field current & observer dependence\\
\((M,P_{\mathrm{Cu}},U^2)\) & three performance levels & local, multi-valued inverse\\
\bottomrule
\end{tabular}
\end{center}

For \((M,\psi_a,\sigma)\), write
\(\psi_a=\vect c^{\mathsf T}\vect r\),
\(\vect r=[i_d,i_{\mathrm{exc}}]^{\mathsf T}\), and
\(\matr R_r=\operatorname{diag}(r_s,r_e)\). The loss-optimal sharing for a
given active flux is
\begin{equation}
 \vect r^\star(\psi_a)=
 \frac{\matr R_r^{-1}\vect c}
 {\vect c^{\mathsf T}\matr R_r^{-1}\vect c}\,\psi_a.
 \label{eq:loss-optimal-sharing}
\end{equation}
Choosing \(\sigma\) along the loss-orthogonal complement makes
\(\sigma=0\) exactly this optimum. Below base speed that is a natural third
axis. In field weakening, a voltage-optimal complement can differ; a smooth
blend is possible, but its frame-rate term must be retained.

\subsection{Performance coordinates}

For the linear EESM let
\[
 M=k_Ti_q(ai_d+bi_e),\quad
 P=r_s(i_d^2+i_q^2)+r_ei_e^2,\quad
 U^2=\vect i^{\mathsf T}\matr Q_U\vect i.
\]
The map \(\gls{sym:taskmap}(\vect i)=(M,P,U^2)\) has Jacobian
\begin{equation}
 \matr J_\Phi=
 \begin{bmatrix}
 k_Tai_q & k_T(ai_d+bi_e) & k_Tbi_q\\
 2r_si_d&2r_si_q&2r_ei_e\\
 2(\matr Q_U\vect i)_d&2(\matr Q_U\vect i)_q&
 2(\matr Q_U\vect i)_e
 \end{bmatrix}.
 \label{eq:performance-jacobian}
\end{equation}
Where \(\det\matr J_\Phi\neq0\), torque, loss, and voltage usage are themselves
local coordinates: their curved, ellipsoidal, and cylindrical level sets all
become coordinate planes. The beauty is static. Globally, quadratic symmetries
and disconnected branches let several currents share the same triple. The map
is singular wherever its three gradients are dependent, including the
zero-torque saddle \(i_q=0,\ ai_d+bi_e=0\); it is also endangered by
\(r_e\to0\), aligned loss/voltage gradients, or removal of a torque mechanism.
Near such sets, the inverse noise gain grows like
\(\|\matr J_\Phi^{-1}\|\).

\subsection{Hyperbolic torque--allocation coordinates}

For either torque sign define
\begin{equation}
 z_\psi=\frac{i_T}{\sqrt2}e^\alpha,\qquad
 z_q=\operatorname{sgn}(M)\frac{i_T}{\sqrt2}e^{-\alpha},\qquad
 M=\frac{\kappa}{2}\operatorname{sgn}(M)i_T^2.
 \label{eq:signed-hyperbolic-chart}
\end{equation}
The inverse on a fixed sign branch is
\[
 i_T=\sqrt{2|z_\psi z_q|},\qquad
 \alpha=\frac12\log\left|\frac{z_\psi}{z_q}\right|,\qquad z_0=z_0.
\]
Differentiation gives
\begin{align}
 \dot z_\psi&=\frac{e^\alpha}{\sqrt2}
  (\dot i_T+i_T\dot\alpha),&
 \dot z_q&=\frac{\operatorname{sgn}(M)e^{-\alpha}}{\sqrt2}
  (\dot i_T-i_T\dot\alpha),\\
 \dot M&=\kappa\operatorname{sgn}(M)i_T\dot i_T.
 \label{eq:hyperbolic-rates-representation}
\end{align}
Thus \(\dot i_T\) changes torque, \(\dot\alpha\) reallocates its two factors at
constant \(i_T\), and \(\dot z_0\) supplies a second torque-neutral direction.
The chart fails at \(i_T=0\) and when either factor vanishes; sign and logarithm
branches also make it unsuitable as a single global startup chart.

\begin{figure}[tb]
\centering
\begin{tikzpicture}[scale=1.0,>=Latex,font=\small]
\draw[->] (-3.1,0)--(3.3,0) node[right]{\(z_\psi\)};
\draw[->] (0,-2.6)--(0,2.8) node[above]{\(z_q\)};
\foreach \c in {.45,.9,1.35}
 {
  \draw[Purple!70,domain=.48:3,samples=80] plot ({\x},{\c/\x});
  \draw[Purple!45,domain=-3:-.48,samples=80] plot ({\x},{\c/\x});
 }
\draw[->,Orange,very thick] (1.2,1.12)--(1.55,.87) node[right]{increase \(\alpha\)};
\draw[->,ForestGreen!70!black,very thick] (1.2,1.12)--(1.65,1.54) node[above]{increase \(i_T\)};
\node[draw,rounded corners,fill=MidnightBlue!7,align=left] at (5.6,.7)
 {\(i_T\): crosses torque sheets\\
  \(\alpha\): moves along one sheet\\
  \(z_0\): second EESM null direction};
\end{tikzpicture}
\caption{Hyperbolic task coordinates separate torque amplitude from reciprocal
sharing of its two multiplicative factors.}
\label{fig:hyperbolic-task}
\end{figure}


The earlier square map \(\zeta=(z_\psi+\mathrm jz_q)^2\) likewise puts torque
in one Cartesian component, but it is two-to-one, singular at the origin, and
requires branch memory. It is an illuminating plot and a poor universal
feedback coordinate. In nonlinear flux maps the same local interpretation
survives whenever the chosen chart Jacobian retains rank, but no fixed
hyperbolic formula is then exact over the full current range.
